Our objective supervises \emph{where} a VLM attends as it generates text, aligning each token's saliency map with the image region its referent occupies. The idea is simple: when the model emits a word that names an annotated entity, the saliency for that word should land on that entity.

\textbf{Token-level grounding signal.}
We assume training captions carry token-level grounding annotations. For a generated caption token $c_t$ that refers to one or more image segments, let $\mathbf{M}_t \in \{0,1\}^{H \times W}$ be the binary mask covering its annotated region(s). From the model's attention we form the saliency map $\mathbf{S}_t$ over the $N$ patch tokens (\Cref{sec:background}), bilinearly upsample it to the annotation resolution $H \times W$, and apply a spatial softmax so that $\mathbf{S}_t$ is a distribution over pixels, $\sum_{ij}\mathbf{S}_{t,ij}=1$. We supervise only the set $\mathcal{T}$ of \emph{grounded} tokens---those annotated with at least one segment---and leave all other tokens (function words, punctuation) to the language modeling loss alone.

\textbf{Alignment loss.}
We compare $\mathbf{S}_t$ against the target distribution that places probability uniformly over the annotated region,
\begin{equation}
\mathbf{T}_t = \mathbf{M}_t \,/\, |\mathbf{M}_t|,
\end{equation}
and penalize their divergence. Our primary objective uses the Kullback--Leibler divergence,
\begin{equation}
\mathcal{L}_{\text{align}}
= \frac{1}{|\mathcal{T}|} \sum_{t \in \mathcal{T}}
\mathrm{KL}\!\left(\mathbf{T}_t \,\|\, \mathbf{S}_t\right)
= \frac{1}{|\mathcal{T}|} \sum_{t \in \mathcal{T}} \sum_{i,j}
\mathbf{T}_{t,ij} \log \frac{\mathbf{T}_{t,ij}}{\mathbf{S}_{t,ij}} .
\label{eq:align}
\end{equation}
Because the target $\mathbf{T}_t$ is fixed, minimizing $\mathrm{KL}(\mathbf{T}_t \,\|\, \mathbf{S}_t)$ is mass-covering: it drives $\mathbf{S}_t$ to put probability everywhere inside the mask and is heavily penalized wherever the mask is occupied but the saliency is near zero, encouraging the map to cover the \emph{whole} referent rather than collapse onto a single patch. As a point of comparison we also implement a squared-error variant, $\mathcal{L}_{\text{align}}^{\text{MSE}} = \frac{1}{|\mathcal{T}|}\sum_{t}\sum_{ij}(\mathbf{S}_{t,ij}-\mathbf{T}_{t,ij})^2$, and use the KL form by default.

\textbf{Training objective.}
We add the alignment term to the standard captioning cross-entropy,
\begin{equation}
\mathcal{L}_{\text{total}} = \mathcal{L}_{\text{CE}} + \lambda \cdot \mathcal{L}_{\text{align}},
\label{eq:total}
\end{equation}
where $\lambda$ trades off fluency against grounding. The term introduces no new parameters and no architectural change, so it drops into a standard VLM fine-tuning pipeline. The vision encoder is kept frozen throughout; which of the remaining components actually absorbs the alignment gradient---the multimodal projector, the language model, or both---is a design choice we treat as an ablation (\Cref{sec:experiment}) and dissect mechanistically in \Cref{sec:analysis}.
